\documentclass[10pt]{article}
\usepackage[margin=20mm]{geometry}
\usepackage{amsmath,amssymb,fontspec,unicode-math,xurl,hyperref}
\IfFontExistsTF{DejaVu Serif}{\setmainfont{DejaVu Serif}}{\setmainfont{Latin Modern Roman}}
\setmathfont{latinmodern-math.otf}
\setlength{\parindent}{0pt}\setlength{\parskip}{6pt}\sloppy
\begin{document}
Abstract: The published formal-verification rule conflates the absence of a reported counterexample with proof, and its contract-plausibility test does not establish that the generated contract specifies the requested repair. These are substantive defects in the advertised formal-guarantee component. Static inspection of the pinned release corroborates the status-handling concern. A separate cost-minimum claim conflicts with the paper’s own comparisons. The benchmark success rates are not disproved. Scientific adjudication remains pending.\par
Author: AI-assisted technical review\par
Date:\par
Title: Technical Verification of PatchPilot\par
\# Paper identity and scope\par
Hongwei Li, Yuheng Tang, Shiqi Wang and Wenbo Guo, PatchPilot: A Cost-Efficient Software Engineering Agent with Early Attempts on Formal Verification, ICML 2025, PMLR 267, pp. 35922–35941 [Retained latex: \textbackslash{}cite\{paper\}]. The reviewed published proceedings PDF has 20 pages. References below use 1-based PDF pages, agreeing with the printed article numbers. The software snapshot is commit `84ec3d3fbc34269078f9f9b62b9d390947344944` [Retained latex: \textbackslash{}cite\{code\}]; its equivalence to historical experimental code is unknown.\par
Inspection covered the five-component workflow, patch-success definition, cost and stability arithmetic, the formal-verification procedure, Appendix A’s ranking rule, implementation settings, case studies and prompts. The complete formal-verification source file was read statically; other components were not exhaustively audited. No author code, model API, container, dataset or trained model was executed.\par
\# Necessary distinctions\par
Equation reference: eq:proof, eq:adequacy. Let \(P(x)\) specify admissible inputs, \(Q(x,f(x))\) a generated postcondition, and \(S(x,f(x))\) the behavior required by the issue. A correctness guarantee needs both \[\begin{aligned}
\text{verification evidence}&\Longrightarrow
 \forall x\,[P(x)\Rightarrow Q(x,f(x))],\\
P(x)\wedge Q(x,f(x))&\Longrightarrow S(x,f(x)).
\end{aligned}\] A proof of a weak contract does not establish the intended repair without the second implication and the appropriate domain. Section 3.1 additionally requires preservation of unaffected functionality. Section 5 (pp. 8–9) generates contracts, calls a contract set plausible when an original function produces an assertion counterexample, and calls a patched instance verified when no patched function produces a counterexample. Page 2 advertises a formal guarantee of patch correctness; p. 9 reports eleven formally verified patches.\par
\# Substantive findings\par
Problem 1 (The negative search result does not establish the formal guarantee).  The Section 5 rule fails to establish the two implications above. The supported assessment is Major for the advertised formal-guarantee component.\par
\#\# Original claim and implementation\par
The paragraph headed Formal Verification, p. 9, was checked in the rendered PDF and extracted text. In , lines 193–196 recognize an assertion by the substring `AssertionError`. Lines 350–392 and 395–416 initialize an assertion flag to false and set it only for that substring. They run time-limited `crosshair check` without requesting explicit confirmation of all conditions. The default timeout is 20 seconds (line 814). The before/after flags enter the success-shaped condition at line 482. This is static source inspection, not author-code execution.\par
CrossHair’s maintainer and documentation distinguish exhaustive confirmation, unknown results and execution errors [Retained latex: \textbackslash{}cite\{chmaint,chdocs\}]. A search can stop with unexamined paths and no counterexample. With a pre-patch assertion already established, either empty post-patch output or an import-error string yields a false assertion flag in the release. An unsuccessful import cannot establish a property of the unexecuted function. The independent diagnostic evaluates this Boolean truth table on explicitly constructed strings; these are not claimed to be actual CrossHair runs or outputs from the eleven published cases. The constructed-status calculation is preserved in .\par
Even if infrastructure errors were discarded manually, an unknown search result would not imply the universal statement in [eq:proof] (\#eq:proof). The current documentation interprets the status semantics; the historical CrossHair version is not pinned by the paper and was not reconstructed.\par
\#\# Complete contract-adequacy counterexample\par
The second gap remains even with a perfect verifier. Consider \[x\in\mathbb Z,\quad P(x)\equiv\mathrm{true},\quad
S(x,y)\equiv y=|x|.\] Let the buggy function be \(b(x)=x\), the generated postcondition be \(Q(x,y)\equiv y\geq0\), and the candidate patch be \(f(x)=0\). All functions are deterministic and total, all inputs and outputs have the declared integer type, and the precondition excludes nothing. No branch removal, approximation or external library is involved.\par
The plausibility gate succeeds because \(Q(-1,b(-1))\) is false. The patched contract holds universally because \(Q(x,0)\) is true for every integer. Yet at \(x=1\) the candidate returns \(0\) instead of \(1\), also damaging previously correct nonnegative-input behavior. Thus both semantic gates in the stated procedure can hold while the intended repair fails. A finite list of checked integers is supplied for reproducibility, but the universal identity \(0\geq0\) and the counterexample \(x=1\) provide the analytic argument. The finite diagnostic output is .\par
Nonnegativity is a necessary property of absolute value, not an unrelated assertion. The described plausibility test fails to establish its sufficiency. This example refutes an inference from that test; it does not claim that this contract or patch appeared in a published trial.\par
\#\# Counterevidence, dependency and supported correction\par
The authors describe an early exploration, acknowledge verifier limitations, and also use empirical patch tests. These facts support a weaker contract-testing interpretation. They do not establish complete specifications or proofs for the eleven cases. The guarantee is a named contribution in the title and introduction, and its support depends directly on the invalid inference. The finding does not require failure of the rest of the patching system.\par
The release defaults to applying the benchmark’s golden patch when no external patch file is supplied (lines 687–690 and 729–732), and its default model is `o4-mini` (line 816); the paper’s formal experiment names o3. These observations raise provenance questions, not evidence that the authors actually evaluated golden patches. The actual command, patches, contracts and full condition statuses are needed.\par
A supported reporting repair separates refuted, unresolved/error and explicitly confirmed contracts; documents verifier assumptions; and validates specification adequacy separately. Unresolved cases must remain unresolved. Strengthening the illustrative postcondition to \(y=|x|\) restores adequacy for this example only. No general repaired pipeline or historical certification was verified here.\par
Problem 2 (The minimum-cost claim conflicts with Table 1).  Page 2 claims the lowest cost among leading compared methods, but Table 1 on p. 6 reports \$0.99 for PatchPilot and \$0.78 for OpenHands on Verified, both labeled Claude-3.5-Sonnet.\par
Consequently, \[0.99>0.78,\qquad (0.99-0.78)/0.78=0.269230769\ldots.\] PatchPilot is about 26.9\% more expensive in that comparison. On Lite, \$0.97 exceeds AutoCodeRover’s \$0.65, although those models differ. The matched-model Verified comparison alone refutes the universal minimum claim. Arithmetic output: .\par
This is a supported Minor comparison defect: the table’s measurements, sub-dollar mean cost and useful accuracy–cost balance remain intact. No underlying API-accounting error is established. The narrower tradeoff claim is supported; actual cost reproduction would require historical usage records.\par
\# Other checks and bounded impact\par
The PatchPilot fractions \(136/300=45.33\%\) and \(268/500=53.60\%\) are consistent. Sample standard deviations for \((30,32,35)\) and \((15,20,21)\) are \(2.5166\) and \(3.2146\), agreeing with Section 4.2. Appendix C.2’s intersection/union ratios are \(25/38=0.6579\) and \(12/25=0.48\). These descriptive checks concern the disclosed selected 45-instance subset, drawn from cases previously solved by both methods. They do not establish a general causal effect of planning style. These consistency calculations are recorded in .\par
Appendix A’s rank \(1_{\mathrm{PoC\ fails}}+n_{\mathrm{failed}}/n_{\mathrm{executed}}\) has the stated broad ordering for positive denominators. Its boundary tie between passing only PoC and passing only all benign tests is resolved by an LLM; strict universal PoC priority would therefore be too strong. No material effect was established, so this is retained as clarification.\par
The measured patching performance is not falsified by the formal-guarantee concern. The eleven verification cases require traceable records. No correction was identified in the inspected README; no exhaustive later-version/erratum search or discovery-priority claim was made.\par
\# Conclusion and state\par
Dimension | State\par
Supported impact | Major: formal guarantee. Minor: cost-minimum comparison.\par
Adjudication | Pending; no human review recorded.\par
Repair verification | Example-level algebra; no pipeline or historical rerun.\par
Author adoption | Not requested; unknown.\par
Remaining work | Historical logs, contract adequacy, patch identities.\par
9 H. Li et al., PatchPilot: A Cost-Efficient Software Engineering Agent with Early Attempts on Formal Verification, PMLR 267 (2025), 35922–35941. https://proceedings.mlr.press/v267/li25cf.html. UCSB ML Security, PatchPilot, commit `84ec3d3fbc34269078f9f9b62b9d390947344944`. https://github.com/ucsb-mlsec/PatchPilot. CrossHair, The Details: Contracts, documentation observed as 0.0.112, accessed 24 October 2026. https://crosshair.readthedocs.io/en/latest/contracts.html. P. Schanely, Is CrossHair a Software Verification Tool?, maintainer discussion 156, opened 11 February 2022; accessed 24 October 2026. https://github.com/pschanely/CrossHair/discussions/156.\par
\# Reproducibility and limitations\par
Original PDF SHA-256: . All 53 manifest-listed staged files matched their hashes. Exact bindings, the source extraction, page image, independent script and actual output reside in . The script uses standard-library arithmetic, imports no author code, and labels its synthetic status strings explicitly.\par
The parent AI reviewer performed and integrated this review; no human assessment is recorded. The user’s selected visible-parent label is 6 Pro, but an independent UI observation was unavailable. Exposed runtime identity is recorded separately in provenance. No model switch, Work task or independent chat was created.\par
The final report PDF was compiled from this TeX using the existing isolated `pdflatex` runtime, without software installation. The parent inspected every rendered report page. Hash-bound compiler, dependency and visual-check facts are recorded in the accompanying .\par
\end{document}
